\footnotesize
\begin{lstlisting}[style=CodeListing]

    #include <iostream>
    #include <chrono>
    #include <cmath>
    #include <vector>
    #include <Windows.h>
    #include <thread>
    
    #include "include/omp.h"
    #define MOD 1000000007 
    
    long long sum(long long x, long long y)
    {
        return (x + y) % MOD;
    }
    long long mult(long long a, long long b)
    {
        return ((a % MOD) * (b % MOD)) % MOD;
    }
    long long binpow(long long x, long long n)
    {
        long long res = 1;
        while (n)
        {
            if (n & 1)
                res = (res * x) % MOD;
            x = (x * x) % MOD;
            n >>= 1;
        }
        return res;
    }
    long long divide(long long a, long long b)
    {
        //return mult(a, binpow(b, phi(MOD) - 1)) % MOD;
        return mult(a, binpow(b, MOD - 2)) % MOD;
    }
    
    void calculate_sum(long long tmax)
    {
        long long taylor = 0;
        long long x = 17, nx = 1, fa = 1;
        #pragma omp parallel for
        for(int i = 0; i < tmax; i++)
        {
            if(i > 0) fa = mult(fa, i);
            nx = mult(nx, x);
            taylor = sum(taylor, divide(nx, fa));
        }
    }
    
    std::vector<long long> v;
    std::vector<long long> table;
    
    const unsigned long long MAX = 1000000;
    const long double stepSize = 0.02d;
    const long long tableStepSize = 10;
    
    int main()
    {
        // freopen("input.txt", "r", stdin);
        // freopen("output_NO_OMP.txt", "w", stdout);
        // freopen("outputOMP.txt", "w", stdout);
    
        MEMORYSTATUSEX status;
        status.dwLength = sizeof(status);
        GlobalMemoryStatusEx(&status);
        unsigned long long totalMemory = status.ullTotalPhys;
        double totalMemoryGB = static_cast<double>(totalMemory) / (1024 * 1024 * 1024);
        std::cout << "Total physical memory: " << totalMemoryGB << " GB" << std::endl;
    
        int numThreads = omp_get_max_threads();
        std::cout << "All threads : " << numThreads << std::endl;
        
        omp_set_num_threads(1);
        std::cout << "1 thread : " << std::endl;
        auto begin = std::chrono::high_resolution_clock::now();
        for(long long tmax = MAX * stepSize; tmax <= MAX; tmax += MAX * stepSize)
        {
            calculate_sum(tmax);
            auto end = std::chrono::high_resolution_clock::now();
            auto duration = std::chrono::duration_cast<std::chrono::microseconds>(end - begin);
            v.push_back(duration.count());
        }
        #pragma omp master
        {
            for(int i = 0; i < v.size(); i += 1)
                std::cout << v[i] << ", ";
        }
        std::cout << std::endl;
        for(size_t i = 0; i < v.size(); i += tableStepSize)
            table.push_back(v[i + 9]);
        v.clear();
        
        // Измерение времени выполнения на всех физических потоках
        omp_set_num_threads(numThreads / 2);
        std::cout << numThreads / 2 << " threads (all physical threads): " << std::endl;
        begin = std::chrono::high_resolution_clock::now();
        for(long long tmax = MAX * stepSize; tmax <= MAX; tmax += MAX * stepSize)
        {
            calculate_sum(tmax);
            auto end = std::chrono::high_resolution_clock::now();
            auto duration = std::chrono::duration_cast<std::chrono::microseconds>(end - begin);
            v.push_back(duration.count());
        }
        #pragma omp master
        {
            for(int i = 0; i < v.size(); i += 1)
                std::cout << v[i] << ", ";
        }
        std::cout << std::endl;
        for(size_t i = 0; i < v.size(); i += tableStepSize)
            table.push_back(v[i + 9]);
        v.clear();
    
        // Измерение времени выполнения на всех логических потоках
        omp_set_num_threads(numThreads);
        std::cout << numThreads << " threads (all logical threads): " << std::endl;
        begin = std::chrono::high_resolution_clock::now();
        for(long long tmax = MAX * stepSize; tmax <= MAX; tmax += MAX * stepSize)
        {
            calculate_sum(tmax);
            auto end = std::chrono::high_resolution_clock::now();
            auto duration = std::chrono::duration_cast<std::chrono::microseconds>(end - begin);
            v.push_back(duration.count());
        }
        #pragma omp master
        {
            for(int i = 0; i < v.size(); i += 1)
                std::cout << v[i] << ", ";
        }
        std::cout << std::endl;
        for(size_t i = 0; i < v.size(); i += tableStepSize)
            table.push_back(v[i + 9]);
        v.clear();
    
        std::cout << "Execution time comparison:" << std::endl;
        std::cout << "-------------------------------------------------------------------------" << std::endl;
        std::cout << "| Threads\t| Time (in milliseconds)\t\t\t\t|" << std::endl;
        std::cout << "-------------------------------------------------------------------------" << std::endl;
        std::cout << "| 1\t\t| ";
        for(size_t i = 0; i < 5; i++)
            std::cout << table[i] << " ";
        std::cout << "\t\t|" << std::endl;
        std::cout << "| Half\t\t| ";
        for(size_t i = 5; i < 10; i++)
            std::cout << table[i] << " ";
        std::cout << "\t\t\t|" << std::endl;
        std::cout << "| All \t\t| ";
        for(size_t i = 10; i < 15; i++)
            std::cout << table[i] << " ";
        std::cout << "\t\t\t|" << std::endl;
        std::cout << "-------------------------------------------------------------------------" << std::endl;
        return 0;
    }

\end{lstlisting}
\normalsize